% GENERATED FILE. Edit canonical lessons, then run npm run generate:books.
% canonical-source-hash: fnv1a64:d7bb66b6cb60cbf2
% canonical-lessons: TE-C235-notu, TE-C235-cillara, TE-C235-jitam, TE-C235-labham, TE-C235-nastam

\chapter{Banknote, Loose change, Salary, Profit, Loss}
\label{ch:te-banknote-loose-change-salary-profit-loss}
\hlchaptermodality{235}

\begin{chapteropening}
\textbf{By the end of this chapter:} \emph{I can say a banknote, loose change, a salary, a profit and a loss.}

Complete the last lesson of chapter 235: I can say a banknote, loose change, a salary, a profit and a loss.
\end{chapteropening}

\section[\texorpdfstring{\te{నోటు} (nōṭu)}{నోటు (nōṭu)}]{\te{నోటు} (nōṭu) — a banknote}

\label{lesson:TE-C235-notu}



\begin{quote}
Before the new one: say the Telugu for a plough, then the Telugu for a sickle.
\end{quote}

\subsection*{You'll want to know: \te{నోటు}}
\textbf{\te{నోటు}} — \emph{nōṭu} — \textquotedblleft{}a banknote\textquotedblright{}.

\begin{grammarlens}[title={What you've built}]
One more word for this chapter.
\end{grammarlens}

\subsection*{Your turn}
\begin{itemize}
\raggedright
  \item \emph{Say it:} \emph{nōṭu}
  \item \emph{Say it:} \emph{nōṭu}, once more
  \item \emph{Say it:} say \emph{nōṭu}
  \item \emph{Recall:} say the Telugu for a plough, then the Telugu for a sickle, then say \emph{nōṭu} again
\end{itemize}

\begin{tcolorbox}[breakable,colback=teal!4,colframe=teal!35!black,title={Before you move on}]
What does \te{నోటు} mean? (\textquotedblleft{}A banknote\textquotedblright{}.) Say it once more.
\end{tcolorbox}
\section[\texorpdfstring{\te{చిల్లర} (cillara)}{చిల్లర (cillara)}]{\te{చిల్లర} (cillara) — loose change}

\label{lesson:TE-C235-cillara}



\begin{quote}
Before the new one: say the Telugu for a sickle, then the Telugu for a banknote.
\end{quote}

\subsection*{You'll want to know: \te{చిల్లర}}
\textbf{\te{చిల్లర}} — \emph{cillara} — \textquotedblleft{}loose change\textquotedblright{}.

\begin{grammarlens}[title={What you've built}]
One more word for this chapter.
\end{grammarlens}

\subsection*{Your turn}
\begin{itemize}
\raggedright
  \item \emph{Say it:} \emph{cillara}
  \item \emph{Say it:} \emph{cillara}, once more
  \item \emph{Say it:} say \emph{cillara}
  \item \emph{Recall:} say the Telugu for a sickle, then the Telugu for a banknote, then say \emph{cillara} again
\end{itemize}

\begin{tcolorbox}[breakable,colback=teal!4,colframe=teal!35!black,title={Before you move on}]
What does \te{చిల్లర} mean? (\textquotedblleft{}Loose change\textquotedblright{}.) Say it once more.
\end{tcolorbox}
\section[\texorpdfstring{\te{జీతం} (jītaṁ)}{జీతం (jītaṁ)}]{\te{జీతం} (jītaṁ) — a salary}

\label{lesson:TE-C235-jitam}



\begin{quote}
Before the new one: say the Telugu for a banknote, then the Telugu for loose change.
\end{quote}

\subsection*{You'll want to know: \te{జీతం}}
\textbf{\te{జీతం}} — \emph{jītaṁ} — \textquotedblleft{}a salary\textquotedblright{}.

\begin{grammarlens}[title={What you've built}]
One more word for this chapter.
\end{grammarlens}

\subsection*{Your turn}
\begin{itemize}
\raggedright
  \item \emph{Say it:} \emph{jītaṁ}
  \item \emph{Say it:} \emph{jītaṁ}, once more
  \item \emph{Say it:} say \emph{jītaṁ}
  \item \emph{Recall:} say the Telugu for a banknote, then the Telugu for loose change, then say \emph{jītaṁ} again
\end{itemize}

\begin{tcolorbox}[breakable,colback=teal!4,colframe=teal!35!black,title={Before you move on}]
What does \te{జీతం} mean? (\textquotedblleft{}A salary\textquotedblright{}.) Say it once more.
\end{tcolorbox}
\section[\texorpdfstring{\te{లాభం} (lābhaṁ)}{లాభం (lābhaṁ)}]{\te{లాభం} (lābhaṁ) — a profit, a gain}

\label{lesson:TE-C235-labham}



\begin{quote}
Before the new one: say the Telugu for loose change, then the Telugu for a salary.
\end{quote}

\subsection*{You'll want to know: \te{లాభం}}
\textbf{\te{లాభం}} — \emph{lābhaṁ} — \textquotedblleft{}a profit, a gain\textquotedblright{}.

\begin{grammarlens}[title={What you've built}]
One more word for this chapter.
\end{grammarlens}

\subsection*{Your turn}
\begin{itemize}
\raggedright
  \item \emph{Say it:} \emph{lābhaṁ}
  \item \emph{Say it:} \emph{lābhaṁ}, once more
  \item \emph{Say it:} say \emph{lābhaṁ}
  \item \emph{Recall:} say the Telugu for loose change, then the Telugu for a salary, then say \emph{lābhaṁ} again
\end{itemize}

\begin{tcolorbox}[breakable,colback=teal!4,colframe=teal!35!black,title={Before you move on}]
What does \te{లాభం} mean? (\textquotedblleft{}A profit, a gain\textquotedblright{}.) Say it once more.
\end{tcolorbox}
\section[\texorpdfstring{\te{నష్టం} (naṣṭaṁ)}{నష్టం (naṣṭaṁ)}]{\te{నష్టం} (naṣṭaṁ) — a loss}

\label{lesson:TE-C235-nastam}



\begin{quote}
Before the new one: say the Telugu for a salary, then the Telugu for a profit.
\end{quote}

\subsection*{You'll want to know: \te{నష్టం}}
\textbf{\te{నష్టం}} — \emph{naṣṭaṁ} — \textquotedblleft{}a loss\textquotedblright{}.

\begin{grammarlens}[title={What you've built}]
One more word for this chapter.
\end{grammarlens}

\subsection*{Your turn}
\begin{itemize}
\raggedright
  \item \emph{Say it:} \emph{naṣṭaṁ}
  \item \emph{Say it:} \emph{naṣṭaṁ}, once more
  \item \emph{Say it:} say \emph{naṣṭaṁ}
  \item \emph{Recall:} say the Telugu for a salary, then the Telugu for a profit, then say \emph{naṣṭaṁ} again
\end{itemize}

\begin{tcolorbox}[breakable,colback=teal!4,colframe=teal!35!black,title={Before you move on}]
What does \te{నష్టం} mean? (\textquotedblleft{}A loss\textquotedblright{}.) Say it once more.
\end{tcolorbox}
